\chapter{Energy-Based Models, Markov Random Fields, and Factor Graphs}
\label{ch:ebm}

Chapter~\ref{ch:statmech} derived the Boltzmann distribution from physics;
this chapter adopts it as a \emph{modelling language} for data. The three
formalisms introduced here --- energy-based models, Markov random fields,
and factor graphs --- are three views of the same object, ordered by
increasing structural explicitness. The chapter is deliberately compact:
every definition it contains is used, verbatim, in the research chapters,
where the noisy Markov chain will be an energy-based model whose factor
graph is a tree.

\section{Energy-based models}\label{sec:ebm-def}

An \emph{energy-based model} (EBM) over $x \in \mathcal{X}$ is a probability
distribution specified by an energy function $E_\theta$:
\begin{equation}
  p_\theta(x) \;=\; \frac{e^{-E_\theta(x)}}{Z(\theta)},
  \qquad
  Z(\theta) = \int_{\mathcal{X}} e^{-E_\theta(x)}\, \dd x ,
  \label{eq:ebm}
\end{equation}
the Boltzmann distribution \eqref{eq:boltzmann} with $\beta$ absorbed into
the energy and the physics replaced by a parametric family
\citep{ackley1985learning, lecun2006tutorial}. The appeal of
\eqref{eq:ebm} is that $E_\theta$ is unconstrained --- any function
defines a valid unnormalised density; the price is the partition function.
Every basic operation hits it:
\begin{itemize}
  \item \emph{Evaluation} of $p_\theta(x)$ requires $Z(\theta)$ ---
  intractable in general (Section~\ref{sec:ising}).
  \item \emph{Sampling} requires MCMC (Chapter~\ref{ch:stochproc}), whose
  mixing time can be exponential in the presence of the metastability
  discussed in Section~\ref{sec:ising}.
  \item \emph{Learning} by maximum likelihood requires the gradient
  \begin{equation}
    -\nabla_\theta \log p_\theta(x)
    = \nabla_\theta E_\theta(x)
    - \E_{p_\theta}\!\big[ \nabla_\theta E_\theta \big],
    \label{eq:ebm-gradient}
  \end{equation}
  whose second term --- the model expectation --- again demands sampling
  from $p_\theta$.
\end{itemize}
There are two classical escape routes from the partition function, and this
thesis walks both. The first is the \emph{score}: since
$\nabla_x \log p_\theta(x) = -\nabla_x E_\theta(x)$, the score is free of
$Z$ --- differentiating in $x$ kills the constant. Score matching
(Chapter~\ref{ch:diffusion}) and diffusion models are built on this
observation. The second is \emph{structure}: when the energy decomposes into
local terms, marginals and even $Z$ itself can become tractable by dynamic
programming --- that route is the rest of this chapter.

\section{Markov random fields}\label{sec:mrf}

\begin{definition}[Markov random field]
Let $G = (V, \mathcal{E})$ be an undirected graph with one random variable
$x_i$ per node. The collection $x = (x_i)_{i \in V}$ is a \emph{Markov
random field} (MRF) if it satisfies the local Markov property: for every
node $i$,
\begin{equation}
  p\big(x_i \,\big|\, x_{V \setminus \{i\}}\big)
  \;=\; p\big(x_i \,\big|\, x_{\partial i}\big),
  \label{eq:local-markov}
\end{equation}
where $\partial i$ is the set of neighbours of $i$ in $G$.
\end{definition}

The fundamental structure theorem connects this conditional-independence
definition to the energy picture: by the Hammersley--Clifford theorem
\citep[proved in][]{besag1974spatial}, a strictly positive density is an MRF
on $G$ if and only if it factorises over the \emph{cliques} of $G$,
\begin{equation}
  p(x) \;=\; \frac{1}{Z} \prod_{C \in \mathrm{cliques}(G)} \psi_C(x_C)
  \;=\; \frac{1}{Z}\, \exp\Big[ -\sum_C E_C(x_C) \Big] .
  \label{eq:hammersley-clifford}
\end{equation}
Conditional independence, graphical locality, and additive local energies
are one and the same thing. The Ising model \eqref{eq:ising} is the textbook
example: cliques are edges, and $E_{(i,j)} = -J_{ij}\sigma_i\sigma_j$.

For \emph{our} purposes the decisive example is one-dimensional. The Markov
chain factorisation \eqref{eq:chain-factorisation},
$p(x) = p_0(x_0) \prod_k M(x_{k+1}|x_k)$, is of the form
\eqref{eq:hammersley-clifford} with $G$ a path graph: a Markov chain is an
MRF on a line, with energy
\begin{equation}
  E(x) \;=\; -\log p_0(x_0) \;-\; \sum_{k=0}^{K-2} \log M(x_{k+1} \,|\, x_k),
  \label{eq:chain-energy}
\end{equation}
a sum of \emph{pairwise, nearest-neighbour} terms. For the Gaussian AR(1)
chain each term is quadratic, so $E$ is a quadratic form whose coupling
matrix connects only neighbours --- read: a \emph{tridiagonal precision
matrix}. Chapter~\ref{ch:gaussian} derives this precision explicitly; the
reader should note that its tridiagonality is not an accident of Gaussian
algebra but the Hammersley--Clifford fingerprint of one-dimensional Markov
structure.

\section{Hidden Markov models: chains observed through noise}
\label{sec:hmm}

The research problem of this thesis is, structurally, a \emph{hidden Markov
model} (HMM) \citep{rabiner1989tutorial}: a latent Markov chain
$(a_0, \dots, a_{K-1})$ evolving by \eqref{eq:ar1}, observed only through
independent noisy measurements --- in our case, one OU channel
\eqref{eq:ou-channel} per frame,
\begin{equation}
  x_k \;=\; a_k\, e^{-t} + \sqrt{\Deltat}\,\varepsilon_k,
  \qquad \varepsilon_k \;\text{i.i.d.}\; \Nd(0,1) .
  \label{eq:hmm-observation}
\end{equation}
The joint law of latents and observations inherits both factorisations:
\begin{equation}
  p(a, x)
  \;=\; \underbrace{p_0(a_0) \prod_{k=0}^{K-2} M(a_{k+1}|a_k)}_{\text{chain prior}}
  \;\cdot\;
  \underbrace{\prod_{k=0}^{K-1} \Nd\!\big(x_k;\, a_k e^{-t}, \Deltat\big)}_{\text{local likelihoods}} .
  \label{eq:hmm-joint}
\end{equation}
Conditioning on the observations gives the \emph{posterior}
$p(a | x) \propto p(a, x)$, which is again an MRF on the line: the
likelihood factors touch one variable each and therefore do not enlarge any
clique. This innocent remark is the structural core of the thesis, for two
reasons. First, as shown in Chapter~\ref{ch:diffusion}, the diffusion score
of the noisy sequence is exactly a posterior expectation under
\eqref{eq:hmm-joint} --- so computing scores \emph{is} HMM inference.
Second, posterior inference on chain-shaped MRFs is solvable in $O(K)$ by
dynamic programming --- the forward--backward algorithm for discrete HMMs
\citep{rabiner1989tutorial}, the Kalman smoother for linear-Gaussian ones
\citep{kalman1960new, rauch1965maximum} --- and both are special cases of
the single algorithm this thesis studies: belief propagation.

\section{Factor graphs}\label{sec:factor-graphs}

MRFs record \emph{which} variables interact but not \emph{how} the joint
factorises (a triangle of pairwise terms and one triple term have the same
graph). The representation that makes factorisation itself graphical --- and
on which message passing is most naturally defined --- is the \emph{factor
graph} \citep{kschischang2001factor}.

\begin{definition}[Factor graph]
Given a factorisation $p(x) = \frac{1}{Z}\prod_{a} \psi_a(x_{\partial a})$,
its factor graph is the bipartite graph with a \emph{variable node} for each
$x_i$, a \emph{factor node} for each $\psi_a$, and an edge $(i, a)$ whenever
$x_i \in x_{\partial a}$ (i.e.\ variable $i$ appears in factor $a$).
\end{definition}

A factor graph is a \emph{tree} when it is connected and cycle-free; trees
are the domain where the sum--product algorithm of Chapter~\ref{ch:bp}
computes every marginal exactly \citep{pearl1988probabilistic,
mezard2009information, bishop2006pattern}.

The posterior \eqref{eq:hmm-joint} of our noisy chain has factor graph:
variable nodes $a_0, \dots, a_{K-1}$; pairwise factor nodes
$M_k = M(a_{k+1}|a_k)$ between consecutive variables; and a unary factor
node $L_k = \Nd(x_k; a_k e^{-t}, \Deltat)$ hanging off each variable (the
observations $x_k$ are fixed numbers, not nodes). Pictorially it is a
\emph{caterpillar}: a spine of $K$ variables alternating with $K-1$
transition factors, plus one likelihood leaf per vertebra. Removing any edge
disconnects it; it contains no cycle; it is a tree. Hence --- recording the
conclusion the research chapters will exploit ---

\begin{proposition}\label{prop:chain-tree}
For any Markov-chain prior and any per-frame observation model, the
posterior $p(a \,|\, x)$ factorises over a tree-shaped factor graph.
Exact marginal inference on it costs $O(K)$ message-passing operations.
\end{proposition}

Note what Proposition~\ref{prop:chain-tree} does \emph{not} say: it does not
say the messages are finitely parametrisable. For discrete chains they are
vectors; for the Gaussian chain they close within the Gaussian family
(Chapter~\ref{ch:bp}); for general continuous chains they are
infinite-dimensional functions, and \emph{that} --- not the graph topology
--- is where approximation enters. Keeping ``tree-exactness of the
algorithm'' separate from ``representability of the messages'' is the
conceptual hinge of this thesis, and the distinction the Gaussian closure
results of Chapter~\ref{ch:bp} are designed to quantify.

\section{What learning adds, and why we go exact}\label{sec:ebm-learning}

A learned generative model of sequences --- an EBM with a neural energy, or
the diffusion models of the next chapter --- must discover the structure
that equations \eqref{eq:chain-energy}--\eqref{eq:hmm-joint} hand us for
free. The modelling bet of this thesis is that the exactly solvable case is
informative about the learned one: if even the \emph{exact} posterior
develops long-range couplings as noise grows (Chapter~\ref{ch:gaussian}
makes this precise), then a learned score faces the same fate; and if exact
inference degrades gracefully under locality constraints
(Chapter~\ref{ch:bp}), then local architectures are a principled bet, not
just a cheap one. The next chapter completes the toolkit by explaining how
modern diffusion models train against these objects at scale.
